\documentclass{article}
\usepackage{graphicx} % Required for inserting images
\usepackage{hyperref}

\title{LaTeX2Typst Documentation}
\author{Gianfranco Baccarella e Vito Di Bari}
\date{April 2026}

\begin{document}

    \maketitle
    \tableofcontents

    \section{Introduction}

    \subsection{Contesto dello sviluppo}
    LaTeX rappresenta un potente linguaggio di markup attualmente utilizzato in numerosi contesti accademici e non per la sua potenza rappresentativa di documenti di testo ordinati secondo regole di formattazione ben definite. Nonostante ciò la sua sintassi e semantica possono risultare complesse per dei neofiti che si approcciano per la prima volta ad un nuovo strumento per la formattazione di testi differente dai classici programmi come Word di Microsoft. Per questo esiste Typst un nuovo linguaggio di markup che ha l'obiettivo di ricoprire le stesse funzioni di LaTeX, ma con una sintassi più semplici e senza la necessità di implementare package per funzioni base come le immagini o i commenti.
    Da questo nasce l'idea di creare un Compilatore che possa prendere in input un file di testo in linguaggio LaTeX e fornire in output un file in linguaggio Typst effettuando quindi una traduzione da un linguaggio di markup ad un altro.

    \subsection{Stato dell'arte}
    Al momento della realizzazione del progetto é stata individuata una repository github con lo stesso obiettivo: convertire dell'input in linguagguio LaTeX in output Typst. La repository, denominata \href{https://github.com/scipenai/tylax}{Tylax} realizzata dall'utente scipenai, di cui é possibile provare il funzionamento al seguente \href{https://convert.silkyai.cn}{link}, é ancora in corso di sviluppo, ma dai risultati pubblicati si nota uno sviluppo avanzato che permette già di tradurre molte strutture sintattiche, compresa una sezione separata esclusiva per la matematica. Nonostante ciò Tylax non supporta al momento la traduzione di differenti comandi come: "maketitle" per la realizzazione del titolo e  alcuni blocchi begin/end come "comment" per creare blocchi di commenti e "flush" per l'allineamento del testo. Inoltre le repository si propone di effettuare una traduzione inversa prendendo in input un file in linguaggio Typst e restituendo quindi un file in linguaggio LaTeX.

    \section{Sturuttura del progetto}
    Il software é stato scritto in linguaggio RUST per poter sfruttare la potenzialità della sua libreria \href{https://docs.rs/pest/2.8.6/pest/}{Pest}, un parser general purpose, con cui é stato possibile definire una grammatica di riferimento ed effettuare un parsing sul testo scritto secondo il markup LaTeX per ottenere un ParseTree, una rappresentazione ad albero della struttura sintattica dell'intero file, basata sulla grammatica da noi definita in modo da scomporre nelle varie componenti il testo.
    La struttura del progetto é abbastanza lineare, prendiamo in input il file LaTeX ed ogni componente dà in input il proprio output alla successiva, i vari passaggi possono essere rappresentati cosi:
    \begin{enumerate}
        \item \textbf{Input}: File Latex + LaTeX Grammar ->
        \item Parser -> Parse Tree ->
        \item Semantic -> Abastract Syntatx Tree ->
        \item CodeGen -> \textbf{Output}: File Typst
    \end{enumerate}

    \subsection{Analisi Sintattica - Parsing}
    Come già inizialmente definito, il software prende in input un file scritto in linguaggio LaTeX (input.tex) e per prima cosa quindi viene effettuato il parsing del testo in modo da poter ottenere il precedentemente citato ParseTree. Per effettuare ciò abbiamo deciso di utilizzare la libreria Pest di Rust, la cui ci ha fornito delle funzioni per potere generare un Parser che potesse seguire determinate regole grammaticali definite nel file \textbf{latex.pest}. La grammatica si basa su definite regole che la stringa dovrà rispettare fino a potersi identificare nella regola più "restrittiva" navigando verso il basso dell'albero, quindi possiamo immaginare anche le regola grammaticali come un albero rovesciato in cui in cima, alla root, vi é presente la regola più generale in questo caso \textbf{file}, che si dirama verso il basso definendo altre regole che possono essere composte dai loro nodi figli. Senza entrare nel dettaglio, in questo documento, della grammatica possiamo fare l'esempio iniziale in cui:
%FARE IL CODE LISTING
    file = { SOI ~ item* ~ EOI }
    item = { block | command | text | newlines | linebreak | comment }
    text = { (!("\\" | NEWLINE | comment) ~ ANY)+ }
% CODE LISTING END
    Come possiamo vedere un file può essere composto da zero o più di un item, la definizione delle regole segue delle espressioni molto simili alle regex, ed a sua volta l'item deve essere composto da uno solo dei diversi elementi definiti. Item definisce quindi una regola generale che identifica ogni elemento presente nel file, successivamente si entra nello specifico di questi elementi come nella regola text che definisce come testo qualsiasi elemento che contenga uno o più caratteri che non inizi con la backslash, altrimenti identificherebbe come comando, e che non rispetti la regola "NEWLINE" o "comment".
    Dopo aver generato un Parser basato sulla nostra grammatica, il sofware effettua quindi il parsing del file di input restituendo in output un Parse Tree in cui suddividiamo ogni elemento del file nella rispettiva regola grammaticale ed organizzati come un albero in cui ogni elemento, definito come Pair, possiede:
    \begin{itemize}
        \item \textbf{rule}, la regola grammaticale rispettata
        \item \textbf{span}, un oggetto di tipo Span, che contiene a sua volta \begin{itemize}
                                                                                   \item \textbf{str}, la stringa di riferimento estratta
                                                                                   \item \textbf{range}, il range occupato dai caratteri
        \end{itemize}
        \item \textbf{inner}, gli elementi figli
    \end{itemize}

    \subsection{Analisi Semantica - AST}
    Dopo aver effettuato il parsing sul testo ed aver ottenuto un ParseTree che ci permette di suddividere il testo in componenti sintattici incapsulando la loro logica, possiamo passare ad un analisi semantica del testo prendendo in input il ParseTree ed ottenere in output una rappresentazione intermedia del codice attraverso l'Abstract Syntax Tree in cui rappresentiamo il testo in una nuova rappresentazione logica ad albero in cui possiamo semplificare la rappresentazione del ParseTree, che risulta essere troppo innestato a causa delle varie regole grammaticali, ottenendo cosi un albero meno folto accomunando differenti regole in un unico enumerativo poiché successivamente rappresentati tramite la stessa funzione; ed inoltre per aggiungere o semplificare informazioni come ad esempio il numero di nuove righe che nel ParseTree vengono rappresentate semplicemente come testo "\\n" differentemente dall'AST in l'oggetto \textbf{NewLineNode} possiede un attributo \textbf{count}. Ad esempio la stringa \textit{\\title{LaTeX To Typst}} che nel ParseTree viene scomposta in 7 nodi ognuno padre dell'altro, differentemente nell'AST é rappresentato con solo 3 nodi.

\end{document}
